% !TEX program = pdflatex
% Source boundary: Piter-authored prior work only. No EDE477 Class 6 science readings were opened or used.
\documentclass[11pt]{article}

\usepackage[margin=0.6in]{geometry}
\usepackage[T1]{fontenc}
\usepackage[utf8]{inputenc}
\usepackage{lmodern}
\usepackage{array}
\usepackage{fancyhdr}
\usepackage{hyperref}
\usepackage{parskip}
\usepackage{ragged2e}
\usepackage{setspace}
\usepackage{needspace}
\usepackage[table]{xcolor}

\definecolor{PrimaryRed}{HTML}{8B1E1E}
\definecolor{DeepNavy}{HTML}{15324A}
\definecolor{SoftBlue}{HTML}{DDEEFF}

\hypersetup{
  colorlinks=true,
  linkcolor=DeepNavy,
  urlcolor=DeepNavy,
  pdftitle={Science Experience Reflection},
  pdfauthor={Piter Zacari Garcia Bautista},
  pdfsubject={EDE477 Pre-Reading Reflective Journal},
  pdfkeywords={science, evidence, access, identity, teaching}
}

\setlength{\parindent}{0pt}
\setlength{\parskip}{0.42em}
\setstretch{1.06}

\newcommand{\reviewsection}[1]{%
  \par\Needspace{4\baselineskip}\addvspace{0.10in}%
  \noindent{\large\bfseries\color{PrimaryRed}#1}\par
  \vspace{0.025in}%
}

\begin{document}

\hypersetup{pageanchor=false}
\begin{titlepage}
\newgeometry{margin=0.6in}
\thispagestyle{empty}

\begin{center}
{\large University of Rochester\\}
{\normalsize Warner School of Education\\[0.15in]}

{\color{PrimaryRed}\rule{\textwidth}{1.4pt}}\\[0.12in]
{\Huge\bfseries Science Experience\\[0.07in]
Reflection}\\[0.15in]
{\huge Evidence, Access, and Identity}\\[0.18in]
{\Large EDE477 --- Pre-Reading Reflective Journal}\\[0.12in]

\vfill

{\color{PrimaryRed}\rule{0.78\textwidth}{0.7pt}}\\[0.20in]
\colorbox{SoftBlue}{%
\parbox{0.88\textwidth}{%
\centering\small\vspace{0.02in}
\textbf{Doing Science Before the Reading}\\[0.05in]
A reflection on lived experience, evidence, access, and the kind of science classroom I hope to build.
}%
}
\end{center}

\vfill

\begin{center}
\renewcommand{\arraystretch}{1.18}
\setlength{\tabcolsep}{4pt}
\begin{tabular}{>{\bfseries}r >{\RaggedRight\arraybackslash}p{5.0in}}
Student & Piter Zacari Garcia Bautista \\
Assignment & \textit{\small Science Experience Reflection---Complete before reading!} \\
Course & EDE477 --- Teaching and Learning in the Content Areas \\
Central Question & \textit{\small What do doing and teaching science mean to me?} \\
Experience Base & \textit{\small Dominican schooling and migration, a BIO/STEM course experience, BIOL550 RNA-seq work, and youth science-camp planning} \\
Sequence Boundary & \textit{\small Completed before consulting the associated EDE477 Science Module readings} \\
Due Date & \textit{\small October 6, 2026, 10:00 AM EDT} \\
\end{tabular}
\end{center}

\vfill

\begin{center}
\begin{minipage}{0.92\textwidth}
\centering\itshape
To me, doing science means asking a question, testing what the evidence can show, and changing my mind when it changes.\\
{\color{PrimaryRed}\rule{4.5in}{0.8pt}}
\end{minipage}
\end{center}

\vfill

\begin{center}
\begin{minipage}{0.95\textwidth}
\small
This journal preserves my thinking before the assigned course reading so I can later revisit how the texts challenged, complicated, or strengthened what I already believed. Its response draws only on my lived experience and prior work; no associated EDE477 Science Module reading was consulted.
\end{minipage}
\end{center}

\restoregeometry
\end{titlepage}

\hypersetup{pageanchor=true}
\pagenumbering{arabic}
\pagestyle{fancy}
\fancyhf{}
\lhead{\small Piter Garcia}
\rhead{\small EDE477 --- Science Experience Reflection}
\cfoot{\small\thepage}
\renewcommand{\headrulewidth}{0.4pt}
\setlength{\headheight}{14pt}

\reviewsection{Experience, Evidence, and Science}

My relationship with science is broader than a traditional science classroom. I began school in the Dominican Republic in Spanish and moved to the United States at eighteen, so I learned science vocabulary, lab expectations, and academic English at the same time as the ideas. To me, doing science means asking a question, testing what the evidence can show, and changing my mind when the evidence changes. It also means asking who gets to ask the question, whose experience counts as evidence, and who is treated as the problem when a design does not fit them.

A difficult experience in a BIO/STEM course made this idea especially real for me. The course design and communication did more than make a challenging class difficult. They affected my health, my trust in the institution, and my ability to participate as a neurodivergent and chronically ill student. When I later wrote about it, I kept returning to one question: if an adult graduate student with technical knowledge, documentation habits, and self-advocacy experience can still be harmed by inaccessible course communication, what happens to children who cannot yet name, document, or resist that harm? The harm did not require malice. It required only a narrow design that assumed one kind of learner. What stayed with me was how a science course can make some students appear capable and others appear deficient. That is a science education issue, because it raises questions about access, language, and who is seen as someone who belongs in science.

Science has also been a place where I felt capable. In BIOL550, my team reanalyzed mouse dorsal root ganglion RNA-sequencing data after nerve injury. I realized that a pattern is evidence but does not explain itself. We had to decide which samples belonged, whether quality checks held, what comparison answered our question, and how far the results supported a claim. Documenting that path let other people check our thinking, and being honest about what the data could not show made the work stronger. Uncertainty felt productive rather than embarrassing. That experience changed how I understand evidence: it is not a verdict handed down by a dataset, but a chain of human decisions that others should be able to follow and question.

This is also why equitable diagnostics matter to me. If a measure or a model fits only some people, the science is incomplete, not the people it misses. In my EDU442 self-portrait, \textit{The Puzzle Plan}, I drew my identities as interlocking pieces of one brain-shaped puzzle: Dominican heritage, a Spanish-first education, neurodivergence, chronic illness, and my research in data science and AI fairness. I used a puzzle to reject a fragmented view of myself, because meaning comes from connection. I see science the same way. Observation, model, language, access, and identity are connected decisions that only make sense together, and a classroom that treats them as separate will keep missing students.

\reviewsection{Teaching and Entering the Classroom}

Teaching science should help students investigate questions that matter in their lives and communities. Planning a youth science-camp investigation of how wheel wear might move material through an environment showed me what that can look like. Learners had to connect an observation to a model and notice the model's limits. A filter can capture material without identifying it, and a possible pathway is not proof that every step has happened. I would want students to separate what they observed, what a source supports, and what they still need to test. Asking where an environmental question comes from, and who is affected by the answer, is part of the science too, not an extra added afterward.

I also think the classroom itself should remove barriers before blaming the learner. I want to teach as a learning scientist: someone who studies the learning environment, notices where access breaks down, and redesigns the system instead of explaining the problem away as a deficit in the student. A multilingual student should be able to reason with home language, drawings, and models before producing polished scientific English. A neurodivergent or chronically ill student should be able to observe, measure, record, talk with a partner, or return after a pause and still make evidence-based claims. A barrier of language, energy, sensory load, or tool design is an access barrier, not a thinking barrier, so rigor stays high while the barriers come down. In my own work I call this the difference between access friction and capability.

My current teaching placement, a high school business and marketing classroom, is not a science class, but it is where I am learning how access works in practice. I have seen how much a predictable routine matters: a visible timer, directions given as what, how, and how long, a modeled example before independent work, and a deck designed with Universal Design for Learning so students can enter the task in more than one way. I have also seen what happens when students help build the expectations. When they drafted their own workplace norms around being on time, prepared, productive, and kind, the rules were theirs rather than something handed to them. My mentor teacher's feedback to keep whole-room awareness, by facing the class and not drifting toward one table while others disengage, taught me that access is also about where the teacher's attention goes. I would carry all of this into science: routines that make the investigation predictable so the thinking can be open.

Choice matters here as well. When students can pick a role, a tool, a language, or a way to show what they know, they gain agency, and their thinking becomes visible in more than one form. Technology can help students collect and organize evidence, but using a tool is not the same as asking a scientific question, so I would judge a tool by whether it widens participation and expression. I would also leave room for curiosity, making, and debugging together. Joy is part of access, because students who feel capable and allowed to participate as themselves will take intellectual risks. Through my \textbf{EQUITAS} lens, I would keep asking whose account of an observation is being believed, and whose is being overlooked.

If I were asked to long-term substitute in a science classroom, my honest first response would be both nervous and excited. I would be nervous because science is larger than my biology and computational background, and I would not pretend to know every answer or every standard immediately. I would be excited because I know how to build learning around questions, evidence, and participation. I would begin by learning the curriculum and the students, making expectations visible, and showing how to ask a better question, check evidence, seek help, and revise an explanation. I want students to see that science belongs to people with many identities, and that changing your mind when the evidence changes is part of doing it, not a failure.

\reviewsection{What I Want to Carry Into the Reading}

I am writing this before the reading so that I can see later what the texts change. Right now I hold a few beliefs I want to test. I believe science learning should begin with real questions and real evidence rather than a list of vocabulary to memorize. I believe students make better claims when they can use more than one language, tool, and way of showing what they know. And I believe equity in science is not only about who is in the room, but about whose questions, explanations, and ways of participating are treated as legitimate science.

I also have honest questions. How much structure do students need before open inquiry becomes useful rather than overwhelming, especially for learners who need predictable routines? How can a teacher judge a student's scientific reasoning without confusing it with fluent English, fast processing, or a neat written product? And how can a science classroom keep rigor high and still let students work at a pace and in a form that protects their health and energy? After the reading, I want to return to this journal and note which of these beliefs were strengthened, which were complicated, and which I now see differently.

\end{document}
